\newpage
\pdfbookmark[1]{\infoname}{info} % Adding the info page among the PDF bookmarks

% English information
\begin{info}
  \begin{abstract}
    Linux command-line skills are an important part of systems-oriented courses at the University of Tartu Institute of Computer Science, including ``Operating Systems'' and ``Computer Security''.
    Earlier thesis work supported these courses by introducing a web-based Linux command-line learning environment and later extending it into a cloud-hosted platform.
    These versions demonstrated the educational value of browser-based Linux practice, but they also revealed shortcomings in maintainability, deployment reproducibility, portability, and cost efficiency.
    The aim of this thesis was to modernize the environment while preserving its educational purpose and core workflow.

    The redesigned platform introduced clearer modular boundaries, a provider-agnostic provisioning layer for user-specific Linux environments, an improved authentication layer, and enhanced multilingual content management.
    Deployment automation was implemented with Helm and Terraform.
    The revised platform was deployed and evaluated on UT HPC, Microsoft Azure, and Amazon Web Services.

    The evaluation, conducted during the Computer Security course in February and March 2026, showed that the revised environment preserved the intended educational workflow.
    Students were able to access the environment, start interactive Linux sessions, and complete the relevant tasks without major issues.
    Compared with the previous Azure deployment, the revised Azure architecture significantly reduced the observed baseline cost of keeping the environment available.
    Among the environments evaluated, UT HPC was the most suitable option for institutional use, while Microsoft Azure had the lowest observed cost among the public cloud deployments.
    The UT HPC version of the revised learning environment is available at\url{https://terminal.cs.ut.ee}.
  \end{abstract}

  \keywords{Linux command line, web-based learning environment, software modernization, cloud computing, cost reduction, Kubernetes}

  \cercs{P175 Informatics, systems theory}
% Codes can be found at: https://www.etis.ee/Portal/Classifiers/Index/26
% Example: P170 Computer science, numerical analysis, systems, control
% Example: P175 Informatics, systems theory
\end{info}


% Estonian information
\begin{otherInfo}{estonian}{Veebipõhise Linuxi käsurea õppekeskkonna kaasajastamine ja kulude vähendamine}
  \begin{abstract}
    Linux ja selle käsurida mängivad olulist rolli Tartu Ülikooli kursustel ``Operatsioonisüsteemid'' ja ``Andmeturve''.
    Nende kursuste praktiliste harjutuste toetamiseks loodi varasemate lõputööde raames veebipõhine Linuxi käsurea õppekeskkond, mida hiljem laiendati pilvepõhiseks platvormiks.
    Käesoleva lõputöö eesmärk oli kujundada olemasolev platvorm ümber jätkusuutlikumaks ja paremini ülekantavaks süsteemiks, säilitades samal ajal selle haridusliku eesmärgi ja põhilise kasutusloogika.

    Töös võeti kasutusele uus arhitektuur, millel olid selgemad moodulipiirid, teenusepakkujast sõltumatu kasutajapõhiste Linuxi keskkondade haldamiskiht ning täiustatud autentimine ja mitmekeelse sisu haldus.
    Keskkondade automatiseeritud ülesseadmiseks kasutati tehnoloogiaid Helm ja Terraform.
    Uuendatud platvorm võeti kasutusele ja testiti keskkondades UT HPC, Microsoft Azure ning Amazon Web Services.

    Testimine, mis viidi läbi 2026.\ aasta veebruaris ja märtsis, näitas, et uuendatud platvorm säilitas esialgse kasutusloogika.
    Tudengid said keskkonda sisse logida, alustada interaktiivseid Linuxi sessioone ja lahendada ettenähtud ülesandeid ilma märkimisväärsete probleemideta.
    Võrreldes varasema Azure'i keskkonnaga vähendas uuendatud Azure'i arhitektuur oluliselt platvormi pideva käigushoidmise baaskulu.
    Testitud keskkondade hulgast osutus UT HPC keskkond sobivaimaks ülikoolipõhiseks kasutuseks ning avalikest pilvekeskkondadest oli madalaima vaadeldud kuluga Microsoft Azure.
    Uuendatud õppekeskkonna UT HPC versioon on kättesaadav aadressil\url{https://terminal.cs.ut.ee}.
  \end{abstract}

  \keywords{Linuxi käsurida, veebipõhine õppekeskkond, tarkvara kaasajastamine, pilveteenused, kulude vähendamine, Kubernetes}

  \cercs{P175 Informaatika, süsteemiteooria}
\end{otherInfo}
